\section{Voltage Control Logic}
\label{sec:ctrl-law}

Light/no-load operation uses hysteretic high/low duty, while loaded operation trims a stored duty command from bulk-voltage error. The hysteretic relationship is
\begin{equation}
 s_k=
 \begin{cases}
  1, & \widehat V_{\mathrm{sb}}\leq V_{\mathrm{on}},\\
  0, & \widehat V_{\mathrm{sb}}\geq V_{\mathrm{release}},\\
  s_{k-1}, & \text{otherwise},
 \end{cases}
 \qquad
 d_k=
 \begin{cases}
  d_{\mathrm{high}}, & s_k=1,\\
  d_{\mathrm{low}}, & s_k=0.
 \end{cases}
 \label{eq:ctrl-hysteresis}
\end{equation}
This decision also accounts for the latest and filtered readings, dwell times, and any demand caused by undervoltage. With the no-load guard active, the operating band is from \(V_{\mathrm{set}}-V_{\mathrm{NLdrop}}\) to \(V_{\mathrm{set}}\) within the threshold limits.

On entry to loaded operation, bounded duty is initialized from \texttt{LDUTY} or the current command. Below \texttt{VON}, a voltage-dependent step increases transfer and sets the reported state to \texttt{LBOOST}. Stored duty is reduced above \texttt{VSET} by more than \SI{0.4}{\volt} and can remain unchanged inside that region.

\texttt{SLEW} sets the base step, scaled by actual elapsed time. The correction is increased for larger voltage droop, using multipliers with defined limits. When the fresh or filtered secondary-voltage reading from PA1 reaches \texttt{VMAX}, the low command is selected immediately and stored duty is reduced. \texttt{LDUTY} sets a floor for the stored value; it is not a measurement of average output duty.

Loaded operation is identified using a threshold on the primary-current reading from PA2 and apparent resistance, with a hold time before the loaded state is released. The default threshold is \SI{0.15}{\ampere} in the configured current scale. A quiet, high-voltage condition must persist for \SI{500}{\milli\second} before release. Since the current signal comes from the primary half bridge, apparent resistance is used to classify load rather than measure external load resistance.

Optional \texttt{LINEFF} adds a duty cap as primary voltage rises. The ratio \(V_{\mathrm{pri,ref}}/\widehat V_{\mathrm{pri}}\) is used within the configured scale bounds. The primary-voltage reading from PA0 therefore limits the high command, while the secondary-voltage reading from PA1 determines the voltage state and error.

\begin{minipage}{\linewidth}
\begin{lstlisting}[language=C,basicstyle=\ttfamily\small,caption={Baseline voltage-control sequence. Fault handling is described in Section~\ref{sec:ctrl-faults}.},label={lst:ctrl-pwr1920-summary}]
read and filter PA0, PA1, PA2 from the circular DMA buffer;
update loaded-state classification;
service soft no-load overvoltage holdoff;
evaluate DMA, enabled current, and enabled OVP faults;
if fault: force PA10 low and return;
if active holdoff: keep BLOW command (default zero) and return;

update no-load escape and undervoltage demand;
if loaded steady is eligible:
    trim the bounded duty memory from PA1 voltage;
    select low duty immediately at VMAX;
else:
    retain the hysteretic high/low state within its band;
    apply the no-load cap or permitted escape demand;

apply optional PA0 line feed-forward and duty bounds;
send the requested duty to the PA10 scheduler with sync polling;
\end{lstlisting}
\end{minipage}
